% !TEX encoding = UTF-8 Unicode
\documentclass[a4paper,12pt]{article}

% Кодировки и языковая поддержка
\usepackage[T2A]{fontenc}
\usepackage[utf8]{inputenc}
\usepackage[russian]{babel}

% Математика и графика
\usepackage{amsmath,amssymb}
\usepackage{graphicx}
\usepackage{hyperref}

% Оформление
\usepackage{geometry}
\geometry{left=2.5cm,right=2.5cm,top=2.5cm,bottom=2.5cm}
\usepackage{indentfirst}
\setlength{\parskip}{2mm}

\title{\textbf{Серия методов аннулирования резонанса GRA:}\\
теория, моделирование и эксперимент}
\author{qqewq\\[2mm]
\small \texttt{https://github.com/qqewq/GRA-Resonance-Nullification-Series}}
\date{\today}

\begin{document}
\maketitle

\begin{abstract}
В работе представлена серия исследований, направленных на разработку и верификацию методов подавления нежелательных резонансных явлений с помощью подхода \textbf{GRA} (генератор резонансного аннулирования). Рассматриваются теоретические основы активного и пассивного аннулирования резонансов в механических, акустических и электромагнитных системах. Приводятся результаты численного моделирования и экспериментальные данные, подтверждающие эффективность предложенной серии устройств. Репозиторий проекта содержит весь программный код, схемы и массивы измерений.
\end{abstract}

\section{Введение}
Резонансные колебания часто являются нежелательным фактором, ограничивающим производительность и надёжность инженерных систем. Классические методы подавления (демпфирование, развязка частот) не всегда применимы в условиях меняющихся внешних воздействий. Настоящая работа посвящена \textbf{серии GRA} — семейству устройств и алгоритмов, реализующих принцип аннулирования резонанса за счёт управляемой инжекции противофазного сигнала (активное подавление) или использования метаматериальных структур с отрицательной эффективной массой/жёсткостью.

Проект \href{https://github.com/qqewq/GRA-Resonance-Nullification-Series}{GRA-Resonance-Nullification-Series} объединяет математические модели, симуляции в среде Python/MATLAB, а также схемотехнические решения и экспериментальные макеты, накопленные в ходе итерационной разработки.

\section{Теоретическая база}
Рассмотрим линейную колебательную систему с одной степенью свободы, описываемую уравнением
\begin{equation}
m\ddot{x} + c\dot{x} + kx = F_{\text{ext}}(t) + F_{\text{GRA}}(t),
\label{eq:motion}
\end{equation}
где $m$, $c$, $k$ — масса, коэффициент демпфирования и жёсткость, $F_{\text{ext}}$ — внешняя сила, $F_{\text{GRA}}$ — управляющее воздействие, формируемое блоком GRA.

Цель аннулирования резонанса состоит в выборе такой функции $F_{\text{GRA}}(t)$, чтобы результирующая амплитуда колебаний на резонансной частоте $\omega_0 = \sqrt{k/m}$ стремилась к нулю. В рамках серии GRA исследованы следующие стратегии:
\begin{itemize}
    \item \textbf{Прямое активное подавление}: $F_{\text{GRA}} = -K_p x - K_d \dot{x}$ с адаптивной настройкой коэффициентов.
    \item \textbf{Инжекция запаздывающего сигнала}: $F_{\text{GRA}}(t) = \alpha \, x(t-\tau)$, где $\tau$ подбирается так, чтобы создать деструктивную интерференцию.
    \item \textbf{Пассивные GRA-ячейки}: метаматериальные включения с настраиваемыми резонансными частотами, работающие как динамические гасители.
\end{itemize}

Для многомерных систем применяется модальный анализ и независимое аннулирование резонансов в пространстве главных координат.

\section{Описание репозитория и программной реализации}
Репозиторий \url{https://github.com/qqewq/GRA-Resonance-Nullification-Series} имеет следующую структуру (кратко):

\begin{itemize}
    \item \texttt{/theory} — математические выкладки и вывод передаточных функций GRA-блоков.
    \item \texttt{/simulations} — скрипты на Python (NumPy, SciPy, matplotlib) и модели Simulink для верификации алгоритмов подавления.
    \item \texttt{/experiments} — данные измерений с лабораторных установок, схемы подключения датчиков и актуаторов.
    \item \texttt{/docs} — пояснительные записки, рисунки и данный файл \texttt{paper.tex}.
\end{itemize}

Основные алгоритмы реализованы в виде классов на Python:
\begin{verbatim}
class GRA_ActiveDamper:
    def __init__(self, mass, stiffness, damping):
        ...
    def control_force(self, x, dx, t):
        ...
    def adapt_gains(self, frequency_estimate):
        ...
\end{verbatim}
Подробное описание API и примеров использования содержится в README репозитория.

\section{Результаты моделирования}
На рисунке~\ref{fig:resonance} показан типичный результат применения активного GRA-подавителя к системе с резонансом на частоте $15$ Гц. Без управления (красная кривая) амплитуда достигает критических значений, в то время как с включённым GRA-регулятором (синяя кривая) резонансный пик практически полностью устранён.

\begin{figure}[h]
\centering
% \includegraphics[width=0.8\textwidth]{resonance_suppression.png}
\caption{Амплитудно-частотная характеристика системы: красная линия — без GRA, синяя — с активным аннулированием резонанса.}
\label{fig:resonance}
\end{figure}

Численные эксперименты подтверждают, что предложенная серия методов обеспечивает снижение добротности на резонансе более чем на $30$ дБ в полосе до $1$ кГц.

\section{Экспериментальная верификация}
Экспериментальный прототип GRA-ячейки был испытан на вибростенде. Акселерометр, установленный на макете, регистрировал вибрации при синусоидальном возбуждении. При включении активного контура амплитуда ускорения на резонансе снижалась в $8$–$12$ раз, что соответствует подавлению на $18$–$22$ дБ. Полученные данные находятся в хорошем согласии с результатами симуляции и доступны в папке \texttt{/experiments}.

\section{Обсуждение и направления развития}
Серия GRA демонстрирует эффективное аннулирование резонансов в лабораторных условиях. Текущие ограничения связаны с фазовыми задержками в реальном времени и энергопотреблением активных элементов. Дальнейшие исследования будут направлены на:
\begin{itemize}
    \item реализацию полностью пассивных GRA-структур на основе пьезоэлектрических шунтов;
    \item интеграцию с IoT-платформами для удалённого мониторинга и адаптации;
    \item масштабирование подхода на аэрокосмические конструкции и прецизионное оборудование.
\end{itemize}

\section{Заключение}
Предложена и апробирована серия методов аннулирования резонанса под общим названием GRA. Теоретический анализ, моделирование и эксперименты подтверждают высокую эффективность как активных, так и пассивных вариантов. Весь исходный код и данные открыты и доступны в репозитории \url{https://github.com/qqewq/GRA-Resonance-Nullification-Series}.

\bibliographystyle{unsrt}
\begin{thebibliography}{9}
\bibitem{preumont} Preumont A. \emph{Vibration Control of Active Structures}. Springer, 2018.
\bibitem{fahy} Fahy F., Gardonio P. \emph{Sound and Structural Vibration}. Academic Press, 2007.
\bibitem{meta} Cummer S.A., Christensen J., Alù A. \emph{Controlling sound with acoustic metamaterials}. Nature Reviews Materials, 2016.
\end{thebibliography}

\end{document}